\documentclass[11pt]{article}
\usepackage{mathblatt}
% Präambel des Setzers (werkzeuge/setzer-praeambel.tex, wortgleich); dazu \gteil (Muster T2U)
\makeatletter
\setlength{\mbpfbe}{0mm}% keine Punkte (pruefung.md § 5: nur beim Prüfstein)
\newlength{\gEin}\setlength{\gEin}{0pt}
\let\gSchrift\relax
\renewcommand{\pfkreuzzeile}[1]{\par\noindent\hangindent3.2em\hangafter1\hspace*{1.6em}$\square$\ #1\par}
\newcommand{\gkreuz}[1]{$\square$\ #1\hspace{2em}}
\newcommand{\gkreuzzeile}[2]{\par\noindent\hangindent2.2em\hangafter1\mbox{#1\,$\square$}\ #2\par}% Buchstabe, Option: eine je Zeile (neben Skizze)
% links, was man ansieht, rechts, was man tut; Spaltengrenze überall an derselben Stelle (gemeinsam.md „Satz“)
% v0.17: Fußeintrag „Nr: Wert“ – kurz in einer Box (nie getrennt), lang umbrechbar mit der Nummer am Anfang
% (vorher \mbox: lange Lösungen liefen über den Rand, AS40 138 pt); robust, weil \fusshilfe in die .aux schreibt
\NewDocumentCommand{\gfuss}{m m}{\sbox0{#1:~#2}\ifdim\wd0<0.45\textwidth\mbox{#1:~#2}\else\mbox{#1:}\nobreak\ #2\fi}
\newsavebox{\gbox}\newlength{\gLinks}
% v0.10 (Lehrer 13.10.): keine starre Halbierung – die linke Spalte ist so breit wie die Figur (+ 6 mm), die
% Antwort steht direkt daneben; ist die Figur breiter als 0,62 Zeile, steht die Antwort darunter
\newcommand{\gzwei}[2]{\par\smallskip\noindent\sbox{\gbox}{#1}%
  \setlength{\gLinks}{\dimexpr\wd\gbox+6mm\relax}%
  \ifdim\gLinks>0.62\linewidth
    \usebox{\gbox}\par\smallskip\noindent\begin{minipage}{\linewidth}#2\end{minipage}%
  \else
    \begin{minipage}[t]{\gLinks}\vspace{0pt}\usebox{\gbox}\end{minipage}%
    \begin{minipage}[t]{\dimexpr\linewidth-\gLinks\relax}\vspace{0pt}\raggedright #2\end{minipage}%
  \fi\par}
% 16.10. (Lehrer, PA10 S. 1): Text links, Koordinatensystem rechts daneben; ist das System breiter als 0,58 Zeile,
% steht es wie bisher unter dem Text
% Passt das System knapp nicht mehr auf die Seite (bis 25 % zu hoch), wird es auf den Rest verkleinert statt
% die Seite halb leer zu lassen
\newlength{\gRest}\newlength{\gHoch}
\newcommand{\gtextks}[2]{\sbox{\gbox}{#1}%
  \setlength{\gHoch}{\dimexpr\ht\gbox+\dp\gbox\relax}%
  \ifdim\pagegoal<\maxdimen\setlength{\gRest}{\dimexpr\pagegoal-\pagetotal-4\baselineskip\relax}%
    \ifdim\gHoch>\gRest\ifdim\gRest>0.75\gHoch
      \sbox{\gbox}{\resizebox*{!}{\gRest}{\usebox{\gbox}}}\fi\fi\fi
  \ifdim\wd\gbox>0.58\linewidth #2\par\smallskip\noindent\gpassend{\usebox{\gbox}}\par
  \else\begin{minipage}[t]{\dimexpr\linewidth-\wd\gbox-5mm\relax}\raggedright #2\end{minipage}\hfill
    \begin{minipage}[t]{\wd\gbox}\vspace{0pt}\usebox{\gbox}\end{minipage}\par\fi}
% v0.9 (12.10., pruefung.md § 5 „Marke mit B, Bedingungen rechts“): drittes Argument von pfaufg = Bedingungen
% der Prüfung (v0.10: \pfohnetr = TR durchgestrichen, „FOR“), klein grau im rechten Rand auf der Grundlinie der Nummer, Zeilen mit \\
\newcommand{\gbedrand}[2]{\ifblank{#2}{}{\rlap{\hspace*{\dimexpr#1+2mm\relax}%
  \parbox[t]{15mm}{\raggedright\scriptsize\leavevmode\color{mbgrau}\strut#2}}}}
% Teil eines Bündels (a, b …): Marke links im Rand, Bedingungen rechts im Rand, auf der Zeile des Teils
\newcommand{\gteilrand}[2]{\llap{\mbpfrandmarke{#1}\hspace*{\mbpfnr}}\gbedrand{\linewidth}{#2}}
% Teil eines Bündels: Buchstabe in eigener Spalte, Inhalt eingerückt (\gEin = Einzug, damit die Spaltengrenze von
% \gzwei an derselben Stelle bleibt)
\newenvironment{gbteil}[3]{\par\smallskip\noindent\gteilrand{#2}{#3}\makebox[\mbpfteil][l]{\textbf{#1}}%
  \begin{minipage}[t]{\dimexpr\linewidth-\mbpfteil\relax}\setlength{\gEin}{\mbpfteil}\raggedright\strut\ignorespaces}%
 {\end{minipage}\par\addvspace{6pt}}% v0.18 (PG80 Nr. 1): Abstand, damit die Linie von a) nicht an b) rutscht
% Grafik höchstens so breit wie die Spalte (Bündel nebeneinander, Koordinatensystem ohne Karo)
\newsavebox{\gpbox}
\newcommand{\gpassend}[1]{\sbox{\gpbox}{#1}\ifdim\wd\gpbox>\linewidth\typeout{GPASSEND: verkleinert \the\wd\gpbox\space auf \the\linewidth}\resizebox{\linewidth}{!}{\usebox{\gpbox}}\else\usebox{\gpbox}\fi}
% Skizzenfeld (v0.10, Lehrer 13.10.): Karo ab dem linken Rand über die Zeile, „Skizze“ klein grau oben links im Feld
% (wie \pfskizzenfeld der Vorlage); der rechte Rand bleibt den Bedingungen
\newcommand{\gskizze}[1]{\par\smallskip\noindent\begin{tikzpicture}
  \clip (0,0) rectangle ({\linewidth-1mm},{#1*10mm});
  \draw[black!16,line width=0.3pt,step=5mm] (0,0) grid ({\linewidth},{#1*10mm});
  \node[anchor=north west,inner sep=2pt,font=\scriptsize,text=mbgrau,fill=white] at (0,{#1*10mm}) {Skizze};
  \end{tikzpicture}\par}
\RenewDocumentEnvironment{pfaufg}{m m m}{%
  \begin{lrbox}{\mbaufgabebox}\begin{minipage}[t]{\linewidth}%
  \gSchrift\noindent\hspace*{\gEin}\mbpfrandmarke{#1}\gbedrand{\linewidth-\gEin}{#3}%
  \begin{minipage}[t]{\mbpfnr}\strut\textbf{#2}\end{minipage}%
  \begin{minipage}[t]{\dimexpr\linewidth-\gEin-\mbpfnr\relax}\raggedright\strut\ignorespaces}%
 {\end{minipage}%
  \end{minipage}\end{lrbox}%
  \setlength{\mbaufgabehoehe}{\dimexpr\ht\mbaufgabebox+\dp\mbaufgabebox\relax}%
  \ifdefempty{\gkopfnach}{\par\addvspace{5pt}\Needspace*{\mbaufgabehoehe}}{%
    % Stufenkopf, der mit dieser (ersten) Aufgabe zusammenbleibt: Platz für beide verlangen, dann den Kopf setzen
    % (12.10., Quadratische: ein Kopf stand allein unten, die Aufgabe mit Koordinatensystem auf der nächsten Seite)
    % v0.17 (AS40 „Terme deuten“ allein unten): hinter dem Kopf steht das Sprungziel (\special, ein Knoten, nach
    % dem die Seite brechen durfte) – darum \nobreak nach dem Kopf; Vorhalt 56 pt = Kopf mit Linie und Abständen
    \par\Needspace*{\dimexpr\mbaufgabehoehe+56pt\relax}\gkopfnach\par\nobreak\gdef\gkopfnach{}\addvspace{5pt}}%
  \noindent\usebox{\mbaufgabebox}\par\addvspace{5pt}}
% v0.18 (plan.md § 10, 15.10. l): Abschnittsbuchstabe A, B, C vor dem Namen im Kopf – als Vorsatz, damit der
% Name im TeX gleich bleibt (Bestellkarte, Sprungziele lesen ihn)
\def\gbstakt{}\newcommand{\gbstvor}[1]{\gdef\gbstakt{#1}}
\NewCommandCopy\gpfstufekopfalt\pfstufekopf
\RenewDocumentCommand{\pfstufekopf}{o m m}{\IfValueTF{#1}{\gpfstufekopfalt[#1]}{\gpfstufekopfalt}%
  {\ifx\gbstakt\empty\else\gbstakt\enspace\fi #2}{#3}\gdef\gbstakt{}}
% v0.18 (16.10. e): Begriff grau, ohne Nummer, ohne Rahmen über der ersten Aufgabe der Stufe
\newcommand{\gbegriff}[1]{\par\nobreak\noindent\hspace*{\mbpfnr}\parbox[t]{\dimexpr\linewidth-\mbpfnr\relax}{\raggedright\small\color{mbgrau}#1}\par\nobreak\addvspace{3pt}}
% v0.18 (16.10. g): Lösungsblatt – Kopf mit kräftigem Strich, Abschnitt klein grau als Zwischenzeile
\newcommand{\glkopf}[1]{\noindent{\large\bfseries Lösungen \textperiodcentered{} #1}\par\vspace{3pt}\hrule height 1.2pt\par\medskip}
\newcommand{\glzwischen}[1]{\par\addvspace{8pt}\Needspace*{4\baselineskip}\noindent{\footnotesize\color{mbgrau}#1}\par\nobreak\vspace{1pt}}
% v0.18 (15.10. o, w): Blattende – Kern ohne Rahmen, dann graue Zeile „Mehr: n Aufgaben“
\newcommand{\gkern}[1]{\par\addvspace{14pt}\Needspace*{6\baselineskip}\noindent{\bfseries Kern}\par\nobreak\smallskip
  {\small\raggedright #1\par}}
\newcommand{\gmehr}[1]{\par\addvspace{10pt}\noindent{\footnotesize\color{mbgrau}Mehr: #1 Aufgaben}\par}
% Stufenkopf vormerken: \gkopfvor{<Kopf>} – die nächste pfaufg setzt ihn, sobald beide zusammen Platz haben
\def\gkopfnach{}
\newcommand{\gkopfvor}[1]{\gdef\gkopfnach{#1}}
% Teilaufgabe mit Skizze links in gzwei: Buchstabe in eigener Spalte, darunter Text und Skizze
\newcommand{\gteil}[2]{\makebox[\mbpfteil][l]{\textbf{#1}}\parbox[t]{\dimexpr0.5\textwidth-\gEin-\mbpfnr-5mm-\mbpfteil\relax}{\raggedright #2}}
% Stufenkopf der Vorlage (mathblatt.sty 2026-10-12: Linie, Abstand, 13 pt, \pfjahre) ohne ihren Vorhalt von
% 14 Zeilen: der Setzer fordert vor dem Kopf 0,16 Seitenhöhe (Kopf + erste Aufgabe), mehr nicht
\setlength{\mbpfstufevorhalt}{0pt}
\makeatother
\newcommand{\karo}[1]{\par\smallskip\noindent\begin{tikzpicture}
  \clip (0,0) rectangle ({\linewidth-1mm},{#1*10mm});
  \draw[black!16,line width=0.3pt,step=5mm] (0,0) grid ({\linewidth},{#1*10mm});
  \end{tikzpicture}\par}
\newcommand{\kk}[1]{$\square$\,#1\hspace{0.9em}}
\newcommand{\linie}[1]{\,{\color{black!45}\rule[-1pt]{#1}{0.4pt}}\,}
\newcommand{\pkt}{\textperiodcentered{}}
% v0.14 (14.10.): Stufenkopf mit fünf Punkten (Option = Jahre für die Bestellkarte) oder „zuletzt 19“
\NewDocumentCommand{\gpunkte}{o m}{#2}
\newcommand{\gzuletzt}[1]{zuletzt #1}
% v0.7 (12.10.): Schülerfrage unter dem Stufenkopf; Ankreuzoption als Bild (Buchstabe und Kästchen darunter)
\newcommand{\gfrage}[1]{\par\nobreak\noindent\hspace*{\mbpfnr}{\small\itshape #1}\par\nobreak\addvspace{2pt}}
\newlength{\gbildbreite}\setlength{\gbildbreite}{0.32\linewidth}% Setzer setzt sie je Reihe (vier Bilder: 0,24)
\newcommand{\gbildoption}[3]{\begin{minipage}[t][#3][b]{\gbildbreite}\centering #2\par\smallskip\mbox{#1\,$\square$}\end{minipage}}
% Lineare (12.10.): Bild mit Buchstabe ohne Kästchen (Diagramm zuordnen); Aussage mit „wahr / falsch“ zum Ankreuzen
\newcommand{\gbildname}[3]{\begin{minipage}[t][#3][b]{\gbildbreite}\centering #2\par\smallskip\mbox{\textbf{#1}}\end{minipage}}
\newcommand{\gwahrfalsch}[2]{\par\noindent\hangindent2.2em\hangafter1\makebox[2.2em][l]{(#1)}#2\hfill\mbox{wahr $\square$\quad falsch $\square$}\par}

\newcommand{\gunter}[2]{\par\addvspace{6pt}\noindent{\bfseries #1\enspace #2}\par\nobreak\vspace{2pt}}
\newcommand{\gsternzeile}{\par\vspace{-4pt}\noindent{\footnotesize\color{mbgrau}* nur FOR}\par\medskip}
% 17.10.: Skizze oben bündig mit der ersten Textzeile (beide Spalten hängen an der Oberkante); der Rechenraum (Karo)
% beginnt direkt unter dem Text, links neben dem Graph, und läuft unter dem Graph über die ganze Zeile weiter
% (Streifen links = auf 5 mm aufgerundete Höhe des Graphen unter dem Text, Raster durchgehend; die Fläche des Streifens
% wird von der Karohöhe unter dem Graph abgezogen, auf 0,5 cm aufgerundet – die Gesamtfläche Rechenraum bleibt mindestens #3 × Zeile); #3 = Zeilen Karo in cm
% unter dem Graph (0 = kein Karo). Kein Verkleinern des Graphen.
\newsavebox{\gtbox}\newlength{\gWl}\newlength{\gDh}
\newcommand{\gkaroohne}[1]{\nointerlineskip\noindent\begin{tikzpicture}
  \clip (0,0) rectangle ({\linewidth-1mm},{#1*10mm});
  \draw[black!16,line width=0.3pt,step=5mm] (0,0) grid ({\linewidth},{#1*10mm});
  \end{tikzpicture}\par}
\newcommand{\gtextkaro}[3]{\sbox{\gbox}{#1}%
  \ifdim\wd\gbox>0.58\linewidth #2\par\smallskip\noindent\gpassend{\usebox{\gbox}}\par\ifnum#3>0 \karo{#3}\fi
  \else
    \setlength{\gWl}{\dimexpr\linewidth-\wd\gbox-5mm\relax}%
    \sbox{\gtbox}{\begin{minipage}[t]{\gWl}\raggedright #2\par\end{minipage}}%
    \setlength{\gDh}{\dimexpr\ht\gbox+\dp\gbox-\ht\gtbox-\dp\gtbox-6pt\relax}%
    \leavevmode
    \ifnum#3>0\ifdim\gDh>0pt
      \pgfmathsetmacro{\gxs}{floor(\gWl/5mm)*5mm}\pgfmathsetmacro{\gys}{ceil(\gDh/5mm)*5mm}%
      \begin{minipage}[t]{\gWl}\usebox{\gtbox}\par\vspace{6pt}\nointerlineskip\noindent
        \begin{tikzpicture}
        \clip (0,0) rectangle (\gxs pt,\gys pt);
        \draw[black!16,line width=0.3pt,step=5mm] (0,0) grid (\gxs pt,\gys pt);
        \end{tikzpicture}\par\end{minipage}%
    \else\usebox{\gtbox}\fi\else\usebox{\gtbox}\fi
    \hfill\begin{minipage}[t]{\wd\gbox}\vspace{-\dimexpr\ht\strutbox+2.4ex\relax}\usebox{\gbox}\end{minipage}\par
    \ifnum#3>0 \ifdim\gDh>0pt
      \pgfmathsetmacro{\gnn}{max(0,ceil(2*(#3-(\gys/28.4528)*(\gxs/\linewidth)))/2)}%
      \ifdim\gnn pt>0pt\gkaroohne{\gnn}\fi
    \else\karo{#3}\fi\fi
  \fi}
\begin{document}

\pfheftstil
\pflaufkopf[P10 \pkt{} Prüfung]{Daten}{PS30 B}\pfseitevon
\pfheftkopf{Daten}{}
\pfstufenzeile{Grafische Darstellungen}
\gsternzeile

% ===================== B Grafische Darstellungen =====================
\gkopfvor{\gbstvor{B}\pfstufekopf{Grafische Darstellungen}{}\gunter{B1}{Diagramme lesen}}
% DRUCK 2022-OS-K4c deckt 2022-OS-K4c
\begin{pfaufg}{nach~P10\\’22}{1.}{}
\gtextkaro{\begin{tikzpicture}[font=\footnotesize,line width=0.6pt]\draw[black!14,line width=0.3pt] (0,0) -- (6.325,0);\draw[black!14,line width=0.3pt] (0,0.364) -- (6.325,0.364);\draw[black!14,line width=0.3pt] (0,0.727) -- (6.325,0.727);\draw[black!14,line width=0.3pt] (0,1.091) -- (6.325,1.091);\draw[black!14,line width=0.3pt] (0,1.455) -- (6.325,1.455);\draw[black!14,line width=0.3pt] (0,1.818) -- (6.325,1.818);\draw[black!14,line width=0.3pt] (0,2.182) -- (6.325,2.182);\draw[black!14,line width=0.3pt] (0,2.545) -- (6.325,2.545);\draw[black!14,line width=0.3pt] (0,2.909) -- (6.325,2.909);\draw[black!14,line width=0.3pt] (0,3.273) -- (6.325,3.273);\draw[black!14,line width=0.3pt] (0,3.636) -- (6.325,3.636);\draw[black!14,line width=0.3pt] (0,4) -- (6.325,4);\draw[black!40,line width=0.3pt] (0,0) -- (6.325,0);\node[left,inner sep=2pt] at (0,0) {390};\draw[black!40,line width=0.3pt] (0,0.364) -- (6.325,0.364);\node[left,inner sep=2pt] at (0,0.364) {400};\draw[black!40,line width=0.3pt] (0,0.727) -- (6.325,0.727);\node[left,inner sep=2pt] at (0,0.727) {410};\draw[black!40,line width=0.3pt] (0,1.091) -- (6.325,1.091);\node[left,inner sep=2pt] at (0,1.091) {420};\draw[black!40,line width=0.3pt] (0,1.455) -- (6.325,1.455);\node[left,inner sep=2pt] at (0,1.455) {430};\draw[black!40,line width=0.3pt] (0,1.818) -- (6.325,1.818);\node[left,inner sep=2pt] at (0,1.818) {440};\draw[black!40,line width=0.3pt] (0,2.182) -- (6.325,2.182);\node[left,inner sep=2pt] at (0,2.182) {450};\draw[black!40,line width=0.3pt] (0,2.545) -- (6.325,2.545);\node[left,inner sep=2pt] at (0,2.545) {460};\draw[black!40,line width=0.3pt] (0,2.909) -- (6.325,2.909);\node[left,inner sep=2pt] at (0,2.909) {470};\draw[black!40,line width=0.3pt] (0,3.273) -- (6.325,3.273);\node[left,inner sep=2pt] at (0,3.273) {480};\draw[black!40,line width=0.3pt] (0,3.636) -- (6.325,3.636);\node[left,inner sep=2pt] at (0,3.636) {490};\draw[black!40,line width=0.3pt] (0,4) -- (6.325,4);\node[left,inner sep=2pt] at (0,4) {500};\filldraw[fill=black!30,draw=black] (0.562,0) rectangle (1.162,2.836);\node[above,inner sep=1.5pt,font=\scriptsize\bfseries] at (0.862,2.836) {468};\node[below,inner sep=3pt] at (0.862,0) {2017};\filldraw[fill=black!30,draw=black] (1.712,0) rectangle (2.312,1.527);\node[above,inner sep=1.5pt,font=\scriptsize\bfseries] at (2.012,1.527) {432};\node[below,inner sep=3pt] at (2.012,0) {2018};\filldraw[fill=black!30,draw=black] (2.862,0) rectangle (3.462,3.455);\node[above,inner sep=1.5pt,font=\scriptsize\bfseries] at (3.162,3.455) {485};\node[below,inner sep=3pt] at (3.162,0) {2019};\filldraw[fill=black!30,draw=black] (4.013,0) rectangle (4.612,2.909);\node[above,inner sep=1.5pt,font=\scriptsize\bfseries] at (4.312,2.909) {470};\node[below,inner sep=3pt] at (4.312,0) {2020};\filldraw[fill=black!30,draw=black] (5.162,0) rectangle (5.762,0.364);\node[above,inner sep=1.5pt,font=\scriptsize\bfseries] at (5.462,0.364) {400};\node[below,inner sep=3pt] at (5.462,0) {2021};\draw[-{Stealth[length=2mm]}] (0,0) -- (0,4.45) node[above right,inner sep=1pt] {Anzahl der Jugendlichen};\draw[-{Stealth[length=2mm]}] (0,0) -- (6.625,0) node[below left,inner sep=2pt,yshift=-9pt] {Jahr};\node[above,font=\small\bfseries] at (3.162,4.75) {Bücher lesen von 2017 bis 2021};\end{tikzpicture}}{In den Jahren 2017 bis 2021 wurden jeweils 1\,200 Jugendliche in Deutschland befragt, wie oft sie Bücher lesen. Das Diagramm zeigt, wie viele der Befragten mehrmals pro Woche Bücher lesen. Paolo sieht auf das Diagramm und behauptet: \emph{„Von Jahr zu Jahr lesen immer weniger Jugendliche Bücher. Von 2018 bis 2021 hat sich die Anzahl der bücherlesenden Jugendlichen um mehr als die Hälfte reduziert.“} Entscheide, ob Paolos Aussage wahr ist. Begründe deine Entscheidung.}{4}
\fusshilfe{\gfuss{1}{falsch}}
\end{pfaufg}\vspace{-6pt}
% DRUCK 2023-OS-K6d deckt 2023-OS-K6d
\begin{pfaufg}{nach~P10\\’23}{2.}{}
\gtextkaro{\begin{tabular}[t]{@{}c@{}}{\footnotesize\setlength{\tabcolsep}{4pt}\renewcommand{\arraystretch}{1.2}\begin{tabular}{|l|r|}\hline \textbf{Sprache} & \begin{tabular}{@{}c@{}}\textbf{Anzahl der Menschen}\\\textbf{mit dieser Muttersprache}\\\textbf{(in Millionen)}\end{tabular}\\\hline Arabisch & 290\\\hline Chinesisch & 1\,300\\\hline Englisch & 500\\\hline Hindi & 525\\\hline Spanisch & 389\\\hline\end{tabular}}\\[2mm]\begin{tikzpicture}[x=5mm,y=5mm,font=\footnotesize,line width=0.6pt]\draw[black!16,line width=0.3pt] (0,0) grid (13,14);\filldraw[fill=black!45,draw=black] (2,0) rectangle (4,7.78);\filldraw[fill=black!45,draw=black] (6,0) rectangle (8,10.5);\foreach \y in {1,...,14}{\draw (-0.15,\y) -- (0.15,\y);}\draw[-{Stealth[length=2mm]}] (0,0) -- (0,14.7) node[above right,inner sep=1pt] {Anzahl in Mio.};\draw[-{Stealth[length=2mm]}] (0,0) -- (13.5,0) node[right,inner sep=2pt] {Sprache};\node[below,inner sep=3pt] at (3,0) {Spanisch};\draw[dashed,line width=0.4pt] (5.6,-0.3) rectangle (8.4,-1.4);\node[below,inner sep=3pt] at (10,0) {Arabisch};\end{tikzpicture}\end{tabular}}{Es gibt vermutlich mehr als 7\,000 Sprachen auf der Welt. In der Tabelle sind die fünf meistgesprochenen Sprachen angegeben. Einige Daten aus der Tabelle sind im Diagramm dargestellt. Ergänze die fehlende Einteilung der $y$-Achse. Notiere unter der zweiten Säule die entsprechende Sprache. Zeichne im Diagramm die Säule für „Arabisch“ ein.}{3}
\fusshilfe{\gfuss{2}{50 Mio.; Hindi; 5,8 Kästchen}}
\end{pfaufg}\vspace{-6pt}
% DRUCK 2026-FOR-K3e deckt 2026-FOR-K3e
\begin{pfaufg}{nach~P10\\’26}{3.\textsuperscript{*}}{}
\gtextkaro{\begin{tabular}[t]{@{}c@{}}{\footnotesize\setlength{\tabcolsep}{4pt}\renewcommand{\arraystretch}{1.2}\begin{tabular}{|l|c|c|c|c|c|c|c|}\hline Wochentag & Mo & Di & Mi & Do & Fr & Sa & So\\\hline \begin{tabular}{@{}l@{}}Benzinpreis in\\€ pro Liter\end{tabular} & 1,77 & 1,78 & 1,77 & 1,79 & 1,84 & 1,82 & 1,85\\\hline\end{tabular}}\\[2mm]\begin{tikzpicture}[font=\footnotesize,line width=0.6pt]\draw[black!14,line width=0.3pt] (0,0) -- (6.375,0);\draw[black!14,line width=0.3pt] (0,0.22) -- (6.375,0.22);\draw[black!14,line width=0.3pt] (0,0.44) -- (6.375,0.44);\draw[black!14,line width=0.3pt] (0,0.66) -- (6.375,0.66);\draw[black!14,line width=0.3pt] (0,0.88) -- (6.375,0.88);\draw[black!14,line width=0.3pt] (0,1.1) -- (6.375,1.1);\draw[black!14,line width=0.3pt] (0,1.32) -- (6.375,1.32);\draw[black!14,line width=0.3pt] (0,1.54) -- (6.375,1.54);\draw[black!14,line width=0.3pt] (0,1.76) -- (6.375,1.76);\draw[black!14,line width=0.3pt] (0,1.98) -- (6.375,1.98);\draw[black!14,line width=0.3pt] (0,2.2) -- (6.375,2.2);\draw[black!14,line width=0.3pt] (0,2.42) -- (6.375,2.42);\draw[black!14,line width=0.3pt] (0,2.64) -- (6.375,2.64);\draw[black!14,line width=0.3pt] (0,2.86) -- (6.375,2.86);\draw[black!14,line width=0.3pt] (0,3.08) -- (6.375,3.08);\draw[black!14,line width=0.3pt] (0,3.3) -- (6.375,3.3);\draw[black!40,line width=0.3pt] (0,0) -- (6.375,0);\node[left,inner sep=2pt] at (0,0) {1,75};\draw[black!40,line width=0.3pt] (0,1.1) -- (6.375,1.1);\node[left,inner sep=2pt] at (0,1.1) {1,80};\draw[black!40,line width=0.3pt] (0,2.2) -- (6.375,2.2);\node[left,inner sep=2pt] at (0,2.2) {1,85};\draw[black!40,line width=0.3pt] (0,3.3) -- (6.375,3.3);\node[left,inner sep=2pt] at (0,3.3) {1,90};\filldraw[fill=black!30,draw=black] (0.412,0) rectangle (0.862,0.44);\node[below,inner sep=3pt] at (0.637,0) {Mo};\filldraw[fill=black!30,draw=black] (1.262,0) rectangle (1.713,0.66);\node[below,inner sep=3pt] at (1.488,0) {Di};\filldraw[fill=black!30,draw=black] (2.112,0) rectangle (2.562,0.44);\node[below,inner sep=3pt] at (2.337,0) {Mi};\filldraw[fill=black!30,draw=black] (2.962,0) rectangle (3.413,0.88);\node[below,inner sep=3pt] at (3.188,0) {Do};\filldraw[fill=black!30,draw=black] (3.812,0) rectangle (4.262,1.98);\node[below,inner sep=3pt] at (4.037,0) {Fr};\filldraw[fill=black!30,draw=black] (4.663,0) rectangle (5.112,1.54);\node[below,inner sep=3pt] at (4.888,0) {Sa};\filldraw[fill=black!30,draw=black] (5.513,0) rectangle (5.962,2.2);\node[below,inner sep=3pt] at (5.737,0) {So};\draw[-{Stealth[length=2mm]}] (0,0) -- (0,3.75) node[above right,inner sep=1pt] {Benzinpreis in € pro Liter};\draw[-{Stealth[length=2mm]}] (0,0) -- (6.675,0) node[below left,inner sep=2pt,yshift=-9pt] {Wochentag};\end{tikzpicture}\end{tabular}}{An einer Tankstelle wurden die Benzinpreise eine Woche lang täglich um 18:00~Uhr notiert. Fabio erstellt aus den Daten dieses Diagramm. Er schreibt in seiner Internetbewertung über die Tankstelle: \emph{„Mein Diagramm zeigt, dass sich die Benzinpreise von Donnerstag auf Freitag mehr als verdoppeln.“} Begründe, warum dieser falsche Eindruck durch Fabios Diagramm entsteht.}{4}
\fusshilfe{\gfuss{3}{Achse ab 1,75 € statt ab 0}}
\end{pfaufg}\vspace{-6pt}
\par\medskip\noindent\textcolor{mbgrau}{\rule{\linewidth}{0.3pt}}\par\nobreak\noindent{\footnotesize\textcolor{mbgrau}{2014–2021}}\par\nobreak\smallskip
% DRUCK 2020-OS-K2b deckt 2020-OS-K2b
\begin{pfaufg}{nach~P10\\’20}{4.}{}
\gtextkaro{\begin{tikzpicture}[font=\footnotesize,line width=0.6pt]\draw[black!14,line width=0.3pt] (0,0) -- (0,3.552);\draw[black!14,line width=0.3pt] (0.22,0) -- (0.22,3.552);\draw[black!14,line width=0.3pt] (0.44,0) -- (0.44,3.552);\draw[black!14,line width=0.3pt] (0.66,0) -- (0.66,3.552);\draw[black!14,line width=0.3pt] (0.88,0) -- (0.88,3.552);\draw[black!14,line width=0.3pt] (1.1,0) -- (1.1,3.552);\draw[black!14,line width=0.3pt] (1.32,0) -- (1.32,3.552);\draw[black!14,line width=0.3pt] (1.54,0) -- (1.54,3.552);\draw[black!14,line width=0.3pt] (1.76,0) -- (1.76,3.552);\draw[black!14,line width=0.3pt] (1.98,0) -- (1.98,3.552);\draw[black!14,line width=0.3pt] (2.2,0) -- (2.2,3.552);\draw[black!14,line width=0.3pt] (2.42,0) -- (2.42,3.552);\draw[black!14,line width=0.3pt] (2.64,0) -- (2.64,3.552);\draw[black!14,line width=0.3pt] (2.86,0) -- (2.86,3.552);\draw[black!14,line width=0.3pt] (3.08,0) -- (3.08,3.552);\draw[black!14,line width=0.3pt] (3.3,0) -- (3.3,3.552);\draw[black!14,line width=0.3pt] (3.52,0) -- (3.52,3.552);\draw[black!14,line width=0.3pt] (3.74,0) -- (3.74,3.552);\draw[black!14,line width=0.3pt] (3.96,0) -- (3.96,3.552);\draw[black!14,line width=0.3pt] (4.18,0) -- (4.18,3.552);\draw[black!14,line width=0.3pt] (4.4,0) -- (4.4,3.552);\draw[black!14,line width=0.3pt] (4.62,0) -- (4.62,3.552);\draw[black!14,line width=0.3pt] (4.84,0) -- (4.84,3.552);\draw[black!14,line width=0.3pt] (5.06,0) -- (5.06,3.552);\draw[black!14,line width=0.3pt] (5.28,0) -- (5.28,3.552);\draw[black!14,line width=0.3pt] (5.5,0) -- (5.5,3.552);\draw[black!14,line width=0.3pt] (5.72,0) -- (5.72,3.552);\draw[black!14,line width=0.3pt] (5.94,0) -- (5.94,3.552);\draw[black!14,line width=0.3pt] (6.16,0) -- (6.16,3.552);\draw[black!14,line width=0.3pt] (6.38,0) -- (6.38,3.552);\draw[black!14,line width=0.3pt] (6.6,0) -- (6.6,3.552);\draw[black!40,line width=0.3pt] (0,0) -- (0,3.552);\node[below,inner sep=2pt] at (0,0) {0};\draw[black!40,line width=0.3pt] (1.1,0) -- (1.1,3.552);\node[below,inner sep=2pt] at (1.1,0) {50};\draw[black!40,line width=0.3pt] (2.2,0) -- (2.2,3.552);\node[below,inner sep=2pt] at (2.2,0) {100};\draw[black!40,line width=0.3pt] (3.3,0) -- (3.3,3.552);\node[below,inner sep=2pt] at (3.3,0) {150};\draw[black!40,line width=0.3pt] (4.4,0) -- (4.4,3.552);\node[below,inner sep=2pt] at (4.4,0) {200};\draw[black!40,line width=0.3pt] (5.5,0) -- (5.5,3.552);\node[below,inner sep=2pt] at (5.5,0) {250};\draw[black!40,line width=0.3pt] (6.6,0) -- (6.6,3.552);\node[below,inner sep=2pt] at (6.6,0) {300};\filldraw[fill=black!30,draw=black] (0,0.196) rectangle (3.74,0.476);\node[left,inner sep=3pt] at (0,0.336) {Mo};\filldraw[fill=black!30,draw=black] (0,0.676) rectangle (2.86,0.956);\node[left,inner sep=3pt] at (0,0.816) {Di};\filldraw[fill=black!30,draw=black] (0,1.156) rectangle (4.62,1.436);\node[left,inner sep=3pt] at (0,1.296) {Mi};\node[left,inner sep=3pt] at (0,1.776) {Do};\filldraw[fill=black!30,draw=black] (0,2.116) rectangle (4.18,2.396);\node[left,inner sep=3pt] at (0,2.256) {Fr};\filldraw[fill=black!30,draw=black] (0,2.596) rectangle (5.28,2.876);\node[left,inner sep=3pt] at (0,2.736) {Sa};\filldraw[fill=black!30,draw=black] (0,3.076) rectangle (6.16,3.356);\node[left,inner sep=3pt] at (0,3.216) {So};\draw[-{Stealth[length=2mm]}] (0,0) -- (7,0);\draw (0,0) -- (0,3.552);\node[below,inner sep=2pt,yshift=-12pt] at (3.3,0) {Anzahl der Besucher};\node[rotate=90,above,inner sep=14pt] at (0,1.776) {Wochentag};\node[above,font=\small\bfseries] at (3.3,3.602) {Besucher im Schwimmbad};\end{tikzpicture}}{Im Diagramm ist die Anzahl der Besucher eines Schwimmbades in der ersten Ferienwoche der Sommerferien dargestellt. Am Donnerstag waren 180 Besucher im Schwimmbad. Notiere, an welchen Tagen mehr als 190 Besucher im Schwimmbad waren.\par\medskip\noindent \linie{60mm}}{0}
\fusshilfe{\gfuss{4}{Mittwoch, Samstag, Sonntag}}
\end{pfaufg}\vspace{-6pt}
% DRUCK 2016-OS-K2a deckt 2016-OS-K2a
% DRUCK 2016-OS-K2b deckt 2016-OS-K2b
\begin{pfaufg}{nach~P10\\’16}{5.}{}
\gtextkaro{\begin{tikzpicture}[font=\footnotesize,line width=0.6pt]\draw[black!14,line width=0.3pt] (0,0) -- (6.84,0);\draw[black!14,line width=0.3pt] (0,0.12) -- (6.84,0.12);\draw[black!14,line width=0.3pt] (0,0.24) -- (6.84,0.24);\draw[black!14,line width=0.3pt] (0,0.36) -- (6.84,0.36);\draw[black!14,line width=0.3pt] (0,0.48) -- (6.84,0.48);\draw[black!14,line width=0.3pt] (0,0.6) -- (6.84,0.6);\draw[black!14,line width=0.3pt] (0,0.72) -- (6.84,0.72);\draw[black!14,line width=0.3pt] (0,0.84) -- (6.84,0.84);\draw[black!14,line width=0.3pt] (0,0.96) -- (6.84,0.96);\draw[black!14,line width=0.3pt] (0,1.08) -- (6.84,1.08);\draw[black!14,line width=0.3pt] (0,1.2) -- (6.84,1.2);\draw[black!14,line width=0.3pt] (0,1.32) -- (6.84,1.32);\draw[black!14,line width=0.3pt] (0,1.44) -- (6.84,1.44);\draw[black!14,line width=0.3pt] (0,1.56) -- (6.84,1.56);\draw[black!14,line width=0.3pt] (0,1.68) -- (6.84,1.68);\draw[black!14,line width=0.3pt] (0,1.8) -- (6.84,1.8);\draw[black!14,line width=0.3pt] (0,1.92) -- (6.84,1.92);\draw[black!14,line width=0.3pt] (0,2.04) -- (6.84,2.04);\draw[black!14,line width=0.3pt] (0,2.16) -- (6.84,2.16);\draw[black!14,line width=0.3pt] (0,2.28) -- (6.84,2.28);\draw[black!14,line width=0.3pt] (0,2.4) -- (6.84,2.4);\draw[black!14,line width=0.3pt] (0,2.52) -- (6.84,2.52);\draw[black!14,line width=0.3pt] (0,2.64) -- (6.84,2.64);\draw[black!14,line width=0.3pt] (0,2.76) -- (6.84,2.76);\draw[black!14,line width=0.3pt] (0,2.88) -- (6.84,2.88);\draw[black!14,line width=0.3pt] (0,3) -- (6.84,3);\draw[black!14,line width=0.3pt] (0,3.12) -- (6.84,3.12);\draw[black!14,line width=0.3pt] (0,3.24) -- (6.84,3.24);\draw[black!14,line width=0.3pt] (0,3.36) -- (6.84,3.36);\draw[black!14,line width=0.3pt] (0,3.48) -- (6.84,3.48);\draw[black!14,line width=0.3pt] (0,3.6) -- (6.84,3.6);\draw[black!14,line width=0.3pt] (0,3.72) -- (6.84,3.72);\draw[black!14,line width=0.3pt] (0,3.84) -- (6.84,3.84);\draw[black!14,line width=0.3pt] (0,3.96) -- (6.84,3.96);\draw[black!14,line width=0.3pt] (0,4.08) -- (6.84,4.08);\draw[black!14,line width=0.3pt] (0,4.2) -- (6.84,4.2);\draw[black!40,line width=0.3pt] (0,0) -- (6.84,0);\node[left,inner sep=2pt] at (0,0) {0};\draw[black!40,line width=0.3pt] (0,0.6) -- (6.84,0.6);\node[left,inner sep=2pt] at (0,0.6) {100};\draw[black!40,line width=0.3pt] (0,1.2) -- (6.84,1.2);\node[left,inner sep=2pt] at (0,1.2) {200};\draw[black!40,line width=0.3pt] (0,1.8) -- (6.84,1.8);\node[left,inner sep=2pt] at (0,1.8) {300};\draw[black!40,line width=0.3pt] (0,2.4) -- (6.84,2.4);\node[left,inner sep=2pt] at (0,2.4) {400};\draw[black!40,line width=0.3pt] (0,3) -- (6.84,3);\node[left,inner sep=2pt] at (0,3) {500};\draw[black!40,line width=0.3pt] (0,3.6) -- (6.84,3.6);\node[left,inner sep=2pt] at (0,3.6) {600};\draw[black!40,line width=0.3pt] (0,4.2) -- (6.84,4.2);\node[left,inner sep=2pt] at (0,4.2) {700};\filldraw[fill=black!45,draw=black] (0.35,0) rectangle (0.73,3.6);\node[below,inner sep=3pt] at (0.54,0) {2007};\filldraw[fill=black!45,draw=black] (1.07,0) rectangle (1.45,2.52);\node[below,inner sep=3pt] at (1.26,0) {2008};\filldraw[fill=black!45,draw=black] (1.79,0) rectangle (2.17,1.92);\node[below,inner sep=3pt] at (1.98,0) {2009};\filldraw[fill=black!45,draw=black] (2.51,0) rectangle (2.89,2.4);\node[below,inner sep=3pt] at (2.7,0) {2010};\filldraw[fill=black!45,draw=black] (3.23,0) rectangle (3.61,1.56);\node[below,inner sep=3pt] at (3.42,0) {2011};\filldraw[fill=black!45,draw=black] (3.95,0) rectangle (4.33,1.56);\node[below,inner sep=3pt] at (4.14,0) {2012};\filldraw[fill=black!45,draw=black] (4.67,0) rectangle (5.05,1.65);\node[below,inner sep=3pt] at (4.86,0) {2013};\filldraw[fill=black!45,draw=black] (5.39,0) rectangle (5.77,1.8);\node[below,inner sep=3pt] at (5.58,0) {2014};\filldraw[fill=black!45,draw=black] (6.11,0) rectangle (6.49,1.5);\node[below,inner sep=3pt] at (6.3,0) {2015};\draw[-{Stealth[length=2mm]}] (0,0) -- (0,4.65) node[above right,inner sep=1pt] {Besucher in Tausend};\draw[-{Stealth[length=2mm]}] (0,0) -- (7.14,0) node[below left,inner sep=2pt,yshift=-9pt] {Jahr};\node[above,font=\small\bfseries] at (3.42,4.95) {Filmpark-Besucher seit 2007};\end{tikzpicture}}{Im Filmpark Babelsberg wird in jedem Jahr die Anzahl der Besucher gezählt.}{0}
\begin{gbteil}{a)}{}{}
Gib ein Jahr an, in dem die Besucherzahl niedriger als 300\,000 war.
\par\medskip\noindent\hfill Jahr $=$\linie{45mm}
\end{gbteil}
\begin{gbteil}{b)}{}{}
Wie hoch war die Besucherzahl im Jahr 2015?
\par\medskip\noindent\hfill Besucherzahl $=$\linie{45mm}
\end{gbteil}
\fusshilfe{\gfuss{5}{a) 2011, 2012, 2013 oder 2015; b) 250 000}}
\end{pfaufg}\vspace{-6pt}
% DRUCK 2018-OS-K3a deckt 2018-OS-K3a
\begin{pfaufg}{nach~P10\\’18}{6.}{}
\gtextkaro{\begin{tabular}[t]{@{}c@{}}{\footnotesize\setlength{\tabcolsep}{4pt}\renewcommand{\arraystretch}{1.2}\begin{tabular}{|l|r|}\hline \textbf{Lieblingsessen} & \textbf{Anzahl}\\\hline Schnitzel & 80\\\hline Pizza & 90\\\hline Kartoffelgerichte & 160\\\hline Gemüsegerichte & 160\\\hline Fischgerichte & 110\\\hline\end{tabular}}\\[2mm]\begin{tabular}{@{}c@{\hspace{4mm}}c@{}}\begin{tabular}[t]{@{}c@{}}\begin{tikzpicture}[font=\footnotesize,line width=0.6pt,x=0.8cm]\draw[black!40,line width=0.3pt] (0,0) -- (0,2.9);\draw[black!40,line width=0.3pt] (0.55,0) -- (0.55,2.9);\draw[black!40,line width=0.3pt] (1.1,0) -- (1.1,2.9);\draw[black!40,line width=0.3pt] (1.65,0) -- (1.65,2.9);\draw[black!40,line width=0.3pt] (2.2,0) -- (2.2,2.9);\draw[black!40,line width=0.3pt] (2.75,0) -- (2.75,2.9);\draw[black!40,line width=0.3pt] (3.3,0) -- (3.3,2.9);\draw[black!40,line width=0.3pt] (3.85,0) -- (3.85,2.9);\draw[black!40,line width=0.3pt] (4.4,0) -- (4.4,2.9);\draw[black!40,line width=0.3pt] (4.95,0) -- (4.95,2.9);\draw[black!40,line width=0.3pt] (5.5,0) -- (5.5,2.9);\filldraw[fill=black!45,draw=black] (0,0.25) rectangle (3.025,0.55);\node[left,inner sep=3pt] at (0,0.4) {A};\filldraw[fill=black!45,draw=black] (0,0.8) rectangle (2.2,1.1);\node[left,inner sep=3pt] at (0,0.95) {B};\filldraw[fill=black!45,draw=black] (0,1.35) rectangle (2.475,1.65);\node[left,inner sep=3pt] at (0,1.5) {C};\filldraw[fill=black!45,draw=black] (0,1.9) rectangle (4.4,2.2);\node[left,inner sep=3pt] at (0,2.05) {D};\filldraw[fill=black!45,draw=black] (0,2.45) rectangle (4.4,2.75);\node[left,inner sep=3pt] at (0,2.6) {E};\draw[-{Stealth[length=2mm]}] (0,0) -- (6.2,0);\draw (0,0) -- (0,2.9);\foreach \k in {1,...,10}{\draw (0.55*\k,-0.08) -- (0.55*\k,0.08);}\node[below left,inner sep=2pt] at (6.2,-0.35) {Anzahl};\node[above,font=\small\bfseries] at (3,2.95) {Lieblingsessen der Jugendlichen};\end{tikzpicture}\end{tabular} & \begin{tabular}[t]{@{}c@{}}\begin{tabular}[t]{@{}l@{}}E: \rule[-0.4ex]{2.4cm}{0.4pt}\\[2.2mm]D: \rule[-0.4ex]{2.4cm}{0.4pt}\\[2.2mm]C: \rule[-0.4ex]{2.4cm}{0.4pt}\\[2.2mm]B: \rule[-0.4ex]{2.4cm}{0.4pt}\\[2.2mm]A: \rule[-0.4ex]{2.4cm}{0.4pt}\end{tabular}\end{tabular}\end{tabular}\end{tabular}}{1000 Jugendliche wurden nach ihren Lieblingsessen befragt. Die Angaben der Tabelle wurden in einem Balkendiagramm dargestellt. Beschrifte die Einteilung an der Achse „Anzahl“ passend. Ordne den Balken die Lieblingsessen zu.}{3}
\fusshilfe{\gfuss{6}{20; A Fischgerichte, B Schnitzel, C Pizza, D und E Kartoffel- und Gemüsegerichte}}
\end{pfaufg}\vspace{-6pt}
% DRUCK 2018-OS-K3d deckt 2018-OS-K3d
\begin{pfaufg}{nach~P10\\’18}{7.\textsuperscript{*}}{}
\gtextkaro{\begin{tabular}[t]{@{}c@{}}\begin{tikzpicture}[font=\footnotesize,line width=0.6pt]\draw[black!14,line width=0.3pt] (0,0) -- (7.875,0);\draw[black!14,line width=0.3pt] (0,0.433) -- (7.875,0.433);\draw[black!14,line width=0.3pt] (0,0.867) -- (7.875,0.867);\draw[black!14,line width=0.3pt] (0,1.3) -- (7.875,1.3);\draw[black!14,line width=0.3pt] (0,1.733) -- (7.875,1.733);\draw[black!14,line width=0.3pt] (0,2.167) -- (7.875,2.167);\draw[black!14,line width=0.3pt] (0,2.6) -- (7.875,2.6);\draw[black!40,line width=0.3pt] (0,0) -- (7.875,0);\draw[black!40,line width=0.3pt] (0,0.433) -- (7.875,0.433);\draw[black!40,line width=0.3pt] (0,0.867) -- (7.875,0.867);\draw[black!40,line width=0.3pt] (0,1.3) -- (7.875,1.3);\draw[black!40,line width=0.3pt] (0,1.733) -- (7.875,1.733);\draw[black!40,line width=0.3pt] (0,2.167) -- (7.875,2.167);\draw[black!40,line width=0.3pt] (0,2.6) -- (7.875,2.6);\draw (-0.08,0) -- (0.08,0); \node[left,inner sep=2pt] at (0,0) {0\,\%};\draw (-0.08,0.433) -- (0.08,0.433); \node[left,inner sep=2pt] at (0,0.433) {10\,\%};\draw (-0.08,0.867) -- (0.08,0.867); \node[left,inner sep=2pt] at (0,0.867) {20\,\%};\draw (-0.08,1.3) -- (0.08,1.3); \node[left,inner sep=2pt] at (0,1.3) {30\,\%};\draw (-0.08,1.733) -- (0.08,1.733); \node[left,inner sep=2pt] at (0,1.733) {40\,\%};\draw (-0.08,2.167) -- (0.08,2.167); \node[left,inner sep=2pt] at (0,2.167) {50\,\%};\draw (-0.08,2.6) -- (0.08,2.6); \node[left,inner sep=2pt] at (0,2.6) {60\,\%};\filldraw[fill=black!45,draw=black] (0.963,0) rectangle (1.663,0.433);\node[below,inner sep=3pt,align=center,text width=1.6cm] at (1.312,0) {nie};\filldraw[fill=black!45,draw=black] (2.712,0) rectangle (3.413,0.65);\node[below,inner sep=3pt,align=center,text width=1.6cm] at (3.062,0) {1-mal bis 3-mal pro Woche};\filldraw[fill=black!45,draw=black] (4.463,0) rectangle (5.162,2.383);\node[below,inner sep=3pt,align=center,text width=1.6cm] at (4.812,0) {4-mal bis 6-mal pro Woche};\filldraw[fill=black!45,draw=black] (6.213,0) rectangle (6.912,0.867);\node[below,inner sep=3pt,align=center,text width=1.6cm] at (6.562,0) {täglich};\draw[-{Stealth[length=2mm]}] (0,0) -- (0,3.05) node[above right,inner sep=1pt] {};\draw[-{Stealth[length=2mm]}] (0,0) -- (8.175,0);\node[above,font=\small\bfseries] at (3.938,3.35) {So oft essen Frauen Fleisch oder Wurst};\end{tikzpicture}\\[2mm]{\small\setlength{\tabcolsep}{5pt}\renewcommand{\arraystretch}{1.25}\begin{tabular}{|p{5.4cm}|c|c|c|}\hline \textbf{Aussage} & \scriptsize\textbf{richtig} & \scriptsize\textbf{falsch} & \scriptsize\begin{tabular}{@{}c@{}}\textbf{nicht ent-}\\\textbf{scheidbar}\end{tabular}\\\hline Weniger als die Hälfte der Frauen essen 4-mal bis 6-mal in der Woche Fleisch oder Wurst. & $\square$ & $\square$ & $\square$\\\hline \multicolumn{4}{|l|}{Begründung:}\\[13mm]\hline Jede 15. Frau isst 1-mal bis 3-mal in der Woche Fleisch oder Wurst. & $\square$ & $\square$ & $\square$\\\hline \multicolumn{4}{|l|}{Begründung:}\\[13mm]\hline\end{tabular}}\end{tabular}}{Frauen wurden nach ihrem Essverhalten befragt. Untersuche die beiden Aussagen. Kreuze an. Gib jeweils eine Begründung an.}{0}
\fusshilfe{\gfuss{7}{(1) falsch; (2) falsch}}
\end{pfaufg}\vspace{-6pt}
% DRUCK 2019-OS-K5c deckt 2019-OS-K5c
\begin{pfaufg}{nach~P10\\’19}{8.\textsuperscript{*}}{}
\gtextkaro{\begin{minipage}[t]{8.4cm}\vspace{0pt}{\large\bfseries Bücherlesen bald out?}\par\smallskip\small \textbf{Berlin} $\blacksquare$ Der Verband der Buchhändler beklagt, dass von Jahr zu Jahr die Anzahl der Jugendlichen, die Bücher lesen, extrem abnimmt. Es wird befürchtet, dass in wenigen Jahren kein Jugendlicher mehr ein Buch liest.\par\medskip\centering \begin{tikzpicture}[font=\footnotesize,line width=0.6pt]\draw[black!14,line width=0.3pt] (0,0) -- (5.5,0);\draw[black!14,line width=0.3pt] (0,0.32) -- (5.5,0.32);\draw[black!14,line width=0.3pt] (0,0.64) -- (5.5,0.64);\draw[black!14,line width=0.3pt] (0,0.96) -- (5.5,0.96);\draw[black!14,line width=0.3pt] (0,1.28) -- (5.5,1.28);\draw[black!14,line width=0.3pt] (0,1.6) -- (5.5,1.6);\draw[black!14,line width=0.3pt] (0,1.92) -- (5.5,1.92);\draw[black!14,line width=0.3pt] (0,2.24) -- (5.5,2.24);\draw[black!14,line width=0.3pt] (0,2.56) -- (5.5,2.56);\draw[black!14,line width=0.3pt] (0,2.88) -- (5.5,2.88);\draw[black!14,line width=0.3pt] (0,3.2) -- (5.5,3.2);\draw[black!14,line width=0.3pt] (0,0) -- (0,3.2);\draw[black!14,line width=0.3pt] (0.32,0) -- (0.32,3.2);\draw[black!14,line width=0.3pt] (0.64,0) -- (0.64,3.2);\draw[black!14,line width=0.3pt] (0.96,0) -- (0.96,3.2);\draw[black!14,line width=0.3pt] (1.28,0) -- (1.28,3.2);\draw[black!14,line width=0.3pt] (1.6,0) -- (1.6,3.2);\draw[black!14,line width=0.3pt] (1.92,0) -- (1.92,3.2);\draw[black!14,line width=0.3pt] (2.24,0) -- (2.24,3.2);\draw[black!14,line width=0.3pt] (2.56,0) -- (2.56,3.2);\draw[black!14,line width=0.3pt] (2.88,0) -- (2.88,3.2);\draw[black!14,line width=0.3pt] (3.2,0) -- (3.2,3.2);\draw[black!14,line width=0.3pt] (3.52,0) -- (3.52,3.2);\draw[black!14,line width=0.3pt] (3.84,0) -- (3.84,3.2);\draw[black!14,line width=0.3pt] (4.16,0) -- (4.16,3.2);\draw[black!14,line width=0.3pt] (4.48,0) -- (4.48,3.2);\draw[black!14,line width=0.3pt] (4.8,0) -- (4.8,3.2);\draw[black!14,line width=0.3pt] (5.12,0) -- (5.12,3.2);\draw[black!14,line width=0.3pt] (5.44,0) -- (5.44,3.2);\draw[black!14,line width=0.3pt] (0,0) -- (5.5,0);\draw[black!14,line width=0.3pt] (0,0.32) -- (5.5,0.32);\draw[black!14,line width=0.3pt] (0,0.64) -- (5.5,0.64);\draw[black!14,line width=0.3pt] (0,0.96) -- (5.5,0.96);\draw[black!14,line width=0.3pt] (0,1.28) -- (5.5,1.28);\draw[black!14,line width=0.3pt] (0,1.6) -- (5.5,1.6);\draw[black!14,line width=0.3pt] (0,1.92) -- (5.5,1.92);\draw[black!14,line width=0.3pt] (0,2.24) -- (5.5,2.24);\draw[black!14,line width=0.3pt] (0,2.56) -- (5.5,2.56);\draw[black!14,line width=0.3pt] (0,2.88) -- (5.5,2.88);\draw[black!14,line width=0.3pt] (0,3.2) -- (5.5,3.2);\draw (-0.08,0) -- (0.08,0); \node[left,inner sep=2pt] at (0,0) {30};\draw (-0.08,1.6) -- (0.08,1.6); \node[left,inner sep=2pt] at (0,1.6) {};\draw (-0.08,3.2) -- (0.08,3.2); \node[left,inner sep=2pt] at (0,3.2) {50};\filldraw[fill=black!45,draw=black] (0.5,0) rectangle (1,2.24);\node[below,inner sep=3pt] at (0.75,0) {2011};\filldraw[fill=black!45,draw=black] (1.5,0) rectangle (2,1.92);\node[below,inner sep=3pt] at (1.75,0) {2012};\filldraw[fill=black!45,draw=black] (2.5,0) rectangle (3,1.6);\node[below,inner sep=3pt] at (2.75,0) {2013};\filldraw[fill=black!45,draw=black] (3.5,0) rectangle (4,1.44);\node[below,inner sep=3pt] at (3.75,0) {2014};\filldraw[fill=black!45,draw=black] (4.5,0) rectangle (5,0.96);\node[below,inner sep=3pt] at (4.75,0) {2015};\draw[-{Stealth[length=2mm]}] (0,0) -- (0,3.65) node[above right,inner sep=1pt] {Angaben in \%};\draw[-{Stealth[length=2mm]}] (0,0) -- (5.8,0) node[below left,inner sep=2pt,yshift=-9pt] {Jahr};\node[above,font=\small\bfseries] at (2.75,3.95) {Bücherlesende Jugendliche};\end{tikzpicture}\end{minipage}}{Im Jahr 2015 wurde folgender Zeitungsartikel veröffentlicht. Schüler einer 10.~Klasse stimmen dem Zeitungsartikel in zwei Aussagen nicht zu. Nenne eine dieser Aussagen und berichtige sie.}{4}
\fusshilfe{\gfuss{8}{„extrem abnimmt“ falsch}}
\end{pfaufg}\vspace{-6pt}
% DRUCK 2014-OS-K3d deckt 2014-OS-K3d
\begin{pfaufg}{nach~P10\\’14}{9.\textsuperscript{*}}{}
\gtextkaro{\begin{tikzpicture}[font=\footnotesize,line width=0.6pt]\draw[black!14,line width=0.3pt] (0,0) -- (4.725,0);\draw[black!14,line width=0.3pt] (0,0.407) -- (4.725,0.407);\draw[black!14,line width=0.3pt] (0,0.815) -- (4.725,0.815);\draw[black!14,line width=0.3pt] (0,1.222) -- (4.725,1.222);\draw[black!14,line width=0.3pt] (0,1.63) -- (4.725,1.63);\draw[black!14,line width=0.3pt] (0,2.037) -- (4.725,2.037);\draw[black!14,line width=0.3pt] (0,2.444) -- (4.725,2.444);\draw[black!14,line width=0.3pt] (0,2.852) -- (4.725,2.852);\draw[black!14,line width=0.3pt] (0,3.259) -- (4.725,3.259);\draw[black!14,line width=0.3pt] (0,3.667) -- (4.725,3.667);\draw[black!14,line width=0.3pt] (0,4.074) -- (4.725,4.074);\draw[black!14,line width=0.3pt] (0,0) -- (0,4.4);\draw[black!14,line width=0.3pt] (0.407,0) -- (0.407,4.4);\draw[black!14,line width=0.3pt] (0.815,0) -- (0.815,4.4);\draw[black!14,line width=0.3pt] (1.222,0) -- (1.222,4.4);\draw[black!14,line width=0.3pt] (1.63,0) -- (1.63,4.4);\draw[black!14,line width=0.3pt] (2.037,0) -- (2.037,4.4);\draw[black!14,line width=0.3pt] (2.444,0) -- (2.444,4.4);\draw[black!14,line width=0.3pt] (2.852,0) -- (2.852,4.4);\draw[black!14,line width=0.3pt] (3.259,0) -- (3.259,4.4);\draw[black!14,line width=0.3pt] (3.667,0) -- (3.667,4.4);\draw[black!14,line width=0.3pt] (4.074,0) -- (4.074,4.4);\draw[black!14,line width=0.3pt] (4.481,0) -- (4.481,4.4);\draw[black!40,line width=0.3pt] (0,0) -- (4.725,0);\draw[black!40,line width=0.3pt] (0,2.037) -- (4.725,2.037);\draw[black!40,line width=0.3pt] (0,4.074) -- (4.725,4.074);\draw (-0.08,0) -- (0.08,0); \node[left,inner sep=2pt] at (0,0) {1060};\draw (-0.08,2.037) -- (0.08,2.037); \node[left,inner sep=2pt] at (0,2.037) {1065};\draw (-0.08,4.074) -- (0.08,4.074); \node[left,inner sep=2pt] at (0,4.074) {1070};\filldraw[fill=black!35,draw=black] (0.838,0) rectangle (1.188,1.67);\node[below,inner sep=3pt,align=center,text width=1.1cm] at (1.013,0) {nach\\einem\\Jahr};\filldraw[fill=black!35,draw=black] (2.038,0) rectangle (2.688,2.811);\node[below,inner sep=3pt,align=center,text width=1.1cm] at (2.363,0) {nach\\zwei\\Jahren};\filldraw[fill=black!35,draw=black] (3.213,0) rectangle (4.213,3.87);\node[below,inner sep=3pt,align=center,text width=1.1cm] at (3.713,0) {nach\\drei\\Jahren};\draw[-{Stealth[length=2mm]}] (0,0) -- (0,4.85) node[above right,inner sep=1pt] {Guthaben in Euro};\draw[-{Stealth[length=2mm]}] (0,0) -- (5.025,0);\node[above,font=\small\bfseries] at (2.363,5.15) {\fbox{\itshape Tolle Zinsen!}};\end{tikzpicture}}{Bei einer Bank erhält man einen Zinssatz von nur 0,25~\% jährlich. Erkläre, wodurch es der Bank trotzdem gelingt, die Guthabenentwicklung in der Grafik günstig aussehen zu lassen.}{4}
\fusshilfe{\gfuss{9}{Achse ab 1060 € statt ab 0}}
\end{pfaufg}\vspace{-6pt}
\gkopfvor{\gunter{B2}{Diagramme zeichnen}}
% DRUCK 2026-FOR-K3d deckt 2026-FOR-K3d 2026-EBR-K4d
\begin{pfaufg}{nach~P10\\’26}{10.}{}
\gtextkaro{\begin{tabular}[t]{@{}c@{}}{\footnotesize\setlength{\tabcolsep}{4pt}\renewcommand{\arraystretch}{1.2}\begin{tabular}{|l|c|c|c|c|c|c|c|}\hline Wochentag & Mo & Di & Mi & Do & Fr & Sa & So\\\hline \begin{tabular}{@{}l@{}}Benzinpreis in\\€ pro Liter\end{tabular} & 1,77 & 1,78 & 1,77 & 1,79 & 1,84 & 1,82 & 1,85\\\hline\end{tabular}}\\[2mm]\begin{tikzpicture}[font=\footnotesize,line width=0.6pt]\draw[black!14,line width=0.3pt] (0,0) -- (6.375,0);\draw[black!14,line width=0.3pt] (0,0.22) -- (6.375,0.22);\draw[black!14,line width=0.3pt] (0,0.44) -- (6.375,0.44);\draw[black!14,line width=0.3pt] (0,0.66) -- (6.375,0.66);\draw[black!14,line width=0.3pt] (0,0.88) -- (6.375,0.88);\draw[black!14,line width=0.3pt] (0,1.1) -- (6.375,1.1);\draw[black!14,line width=0.3pt] (0,1.32) -- (6.375,1.32);\draw[black!14,line width=0.3pt] (0,1.54) -- (6.375,1.54);\draw[black!14,line width=0.3pt] (0,1.76) -- (6.375,1.76);\draw[black!14,line width=0.3pt] (0,1.98) -- (6.375,1.98);\draw[black!14,line width=0.3pt] (0,2.2) -- (6.375,2.2);\draw[black!14,line width=0.3pt] (0,2.42) -- (6.375,2.42);\draw[black!14,line width=0.3pt] (0,2.64) -- (6.375,2.64);\draw[black!14,line width=0.3pt] (0,2.86) -- (6.375,2.86);\draw[black!14,line width=0.3pt] (0,3.08) -- (6.375,3.08);\draw[black!14,line width=0.3pt] (0,3.3) -- (6.375,3.3);\draw[black!40,line width=0.3pt] (0,0) -- (6.375,0);\node[left,inner sep=2pt] at (0,0) {1,75};\draw[black!40,line width=0.3pt] (0,1.1) -- (6.375,1.1);\node[left,inner sep=2pt] at (0,1.1) {1,80};\draw[black!40,line width=0.3pt] (0,2.2) -- (6.375,2.2);\node[left,inner sep=2pt] at (0,2.2) {1,85};\draw[black!40,line width=0.3pt] (0,3.3) -- (6.375,3.3);\node[left,inner sep=2pt] at (0,3.3) {1,90};\filldraw[fill=black!30,draw=black] (0.412,0) rectangle (0.862,0.44);\node[below,inner sep=3pt] at (0.637,0) {Mo};\filldraw[fill=black!30,draw=black] (1.262,0) rectangle (1.713,0.66);\node[below,inner sep=3pt] at (1.488,0) {Di};\filldraw[fill=black!30,draw=black] (2.112,0) rectangle (2.562,0.44);\node[below,inner sep=3pt] at (2.337,0) {Mi};\filldraw[fill=black!30,draw=black] (2.962,0) rectangle (3.413,0.88);\node[below,inner sep=3pt] at (3.188,0) {Do};\filldraw[fill=black!30,draw=black] (3.812,0) rectangle (4.262,1.98);\node[below,inner sep=3pt] at (4.037,0) {Fr};\filldraw[fill=black!30,draw=black] (4.663,0) rectangle (5.112,1.54);\node[below,inner sep=3pt] at (4.888,0) {Sa};\node[below,inner sep=3pt] at (5.737,0) {So};\draw[-{Stealth[length=2mm]}] (0,0) -- (0,3.75) node[above right,inner sep=1pt] {Benzinpreis in € pro Liter};\draw[-{Stealth[length=2mm]}] (0,0) -- (6.675,0) node[below left,inner sep=2pt,yshift=-9pt] {Wochentag};\end{tikzpicture}\end{tabular}}{An einer Tankstelle wurden die Benzinpreise eine Woche lang täglich um 18:00~Uhr notiert. Fabio erstellt aus den Daten der Tabelle dieses Diagramm. Zeichne die fehlende Säule im Diagramm ein.}{0}
\fusshilfe{\gfuss{10}{Säule So bis 1,85 €}}
\end{pfaufg}\vspace{-6pt}
% DRUCK 2025-OS-K6b deckt 2025-OS-K6b
\begin{pfaufg}{nach~P10\\’25}{11.}{}
\gtextkaro{\begin{tikzpicture}[line width=0.6pt]\draw (0,0) circle (1.8);\fill (0,0) circle (1.4pt);\end{tikzpicture}}{Die Liste zeigt das Alter der 11 Erziehenden in einer Kita: 52 \quad 57 \quad 33 \quad 44 \quad 31 \quad 30 \quad 66 \quad 58 \quad 61 \quad 41 \quad 55. Ermittle den prozentualen Anteil der Erziehenden, die jünger als 50 Jahre sind. Stelle diesen prozentualen Anteil im Kreisdiagramm dar.}{5}
\par\medskip\noindent\hfill Anteil $=$\linie{45mm}
\fusshilfe{\gfuss{11}{45,5 \%; 163,6°}}
\end{pfaufg}\vspace{-6pt}
% DRUCK 2022-OS-K4d deckt 2022-OS-K4d
\begin{pfaufg}{nach~P10\\’22}{12.}{}
\gtextkaro{\begin{tabular}[t]{@{}c@{}}{\footnotesize\setlength{\tabcolsep}{4pt}\renewcommand{\arraystretch}{1.2}\begin{tabular}{|l|r|}\hline mehrmals pro Woche & 39\,\%\\\hline einmal pro Woche & 18\,\%\\\hline einmal im Monat & 27\,\%\\\hline nie & 16\,\%\\\hline\end{tabular}}\\[2mm]\begin{tikzpicture}[font=\footnotesize,line width=0.6pt]\filldraw[fill=white,draw=black,line width=0.5pt] (0,0) -- (90:1.5) arc (90:32.4:1.5) -- cycle;\filldraw[fill=white,draw=black,line width=0.5pt] (0,0) -- (32.4:1.5) arc (32.4:-64.8:1.5) -- cycle;\filldraw[fill=white,draw=black,line width=0.5pt] (0,0) -- (-64.8:1.5) arc (-64.8:-129.6:1.5) -- cycle;\filldraw[fill=white,draw=black,line width=0.5pt] (0,0) -- (-129.6:1.5) arc (-129.6:-270:1.5) -- cycle;\draw (0,0) circle (1.5);\draw[line width=0.4pt] (0.434,0.789) -- (2,1.55) -- (4,1.55);\node[anchor=south east,inner sep=2pt] at (4,1.55) {nie};\draw[line width=0.4pt] (0.864,-0.251) -- (2,-0.2) -- (4,-0.2);\draw[line width=0.4pt] (-0.113,-0.893) -- (2,-1.75) -- (4,-1.75);\draw[line width=0.4pt] (-0.847,0.305) -- (-2,0.2) -- (-4,0.2);\end{tikzpicture}\end{tabular}}{750 Jugendliche wurden befragt, wie häufig sie Comic-Hefte lesen. Die Tabelle zeigt die prozentuale Verteilung. Beschrifte das Kreisdiagramm entsprechend der Tabelle. Berechne die Winkelgröße für den Anteil „nie“.}{3}
\par\medskip\noindent\hfill Winkel $=$\linie{45mm}
\fusshilfe{\gfuss{12}{rechts einmal im Monat, unten einmal pro Woche, links mehrmals pro Woche; 57,6°}}
\end{pfaufg}\vspace{-6pt}
% DRUCK 2024-OS-K2c deckt 2024-OS-K2c
\begin{pfaufg}{nach~P10\\’24}{13.\textsuperscript{*}}{}
\gtextkaro{\begin{tabular}[t]{@{}c@{}}{\footnotesize\setlength{\tabcolsep}{4pt}\renewcommand{\arraystretch}{1.2}\begin{tabular}{|r|l|r|}\hline 1 & Trinken und Kochen & 12 $\ell$\\\hline 2 & Körperpflege & 256 $\ell$\\\hline 3 & Toilettenspülung & 136 $\ell$\\\hline 4 & Wäsche waschen & 60 $\ell$\\\hline 5 & Gartenpflege & 16 $\ell$\\\hline 6 & Sonstiges & 28 $\ell$\\\hline\end{tabular}}\\[2mm]\begin{tikzpicture}[font=\footnotesize,line width=0.6pt]\filldraw[fill=white,draw=black,line width=0.5pt] (0,0) -- (90:1.5) arc (90:-91.417:1.5) -- cycle;\filldraw[fill=white,draw=black,line width=0.5pt] (0,0) -- (-91.417:1.5) arc (-91.417:-111.26:1.5) -- cycle;\filldraw[fill=white,draw=black,line width=0.5pt] (0,0) -- (-111.26:1.5) arc (-111.26:-153.78:1.5) -- cycle;\filldraw[fill=black!25,draw=black,line width=0.5pt] (0,0) -- (-153.78:1.5) arc (-153.78:-165.118:1.5) -- cycle;\filldraw[fill=white,draw=black,line width=0.5pt] (0,0) -- (-165.118:1.5) arc (-165.118:-261.496:1.5) -- cycle;\filldraw[fill=black!25,draw=black,line width=0.5pt] (0,0) -- (-261.496:1.5) arc (-261.496:-270:1.5) -- cycle;\draw (0,0) circle (1.5);\draw[line width=0.4pt] (-0.067,0.898) -- (2,1.95) -- (4,1.95);\node[anchor=south east,inner sep=2pt] at (4,1.95) {1 Trinken und Kochen};\draw[line width=0.4pt] (-0.843,-0.316) -- (-2,0.9) -- (-4,0.9);\node[anchor=south west,inner sep=2pt] at (-4,0.9) {5 Gartenpflege};\end{tikzpicture}\end{tabular}}{Familie Schulz verbraucht täglich ca. 508 Liter Leitungswasser. Die Tabelle zeigt ihren Verbrauch. Im Kreisdiagramm sind die Anteile aus der Tabelle dargestellt. Beschrifte den Anteil für „Wäsche waschen“ im Kreisdiagramm. Berechne die Winkelgröße für den Anteil 1 „Trinken und Kochen“ vom gesamten Wasserverbrauch.}{4}
\par\medskip\noindent\hfill Winkel $=$\linie{45mm}
\fusshilfe{\gfuss{13}{Sektor unter 5 Gartenpflege; 8,5°}}
\end{pfaufg}\vspace{-6pt}
\par\medskip\noindent\textcolor{mbgrau}{\rule{\linewidth}{0.3pt}}\par\nobreak\noindent{\footnotesize\textcolor{mbgrau}{2015–2021}}\par\nobreak\smallskip
% DRUCK 2020-OS-K2a deckt 2020-OS-K2a
\begin{pfaufg}{nach~P10\\’20}{14.}{}
\gtextkaro{\begin{tikzpicture}[font=\footnotesize,line width=0.6pt]\draw[black!14,line width=0.3pt] (0,0) -- (0,3.552);\draw[black!14,line width=0.3pt] (0.22,0) -- (0.22,3.552);\draw[black!14,line width=0.3pt] (0.44,0) -- (0.44,3.552);\draw[black!14,line width=0.3pt] (0.66,0) -- (0.66,3.552);\draw[black!14,line width=0.3pt] (0.88,0) -- (0.88,3.552);\draw[black!14,line width=0.3pt] (1.1,0) -- (1.1,3.552);\draw[black!14,line width=0.3pt] (1.32,0) -- (1.32,3.552);\draw[black!14,line width=0.3pt] (1.54,0) -- (1.54,3.552);\draw[black!14,line width=0.3pt] (1.76,0) -- (1.76,3.552);\draw[black!14,line width=0.3pt] (1.98,0) -- (1.98,3.552);\draw[black!14,line width=0.3pt] (2.2,0) -- (2.2,3.552);\draw[black!14,line width=0.3pt] (2.42,0) -- (2.42,3.552);\draw[black!14,line width=0.3pt] (2.64,0) -- (2.64,3.552);\draw[black!14,line width=0.3pt] (2.86,0) -- (2.86,3.552);\draw[black!14,line width=0.3pt] (3.08,0) -- (3.08,3.552);\draw[black!14,line width=0.3pt] (3.3,0) -- (3.3,3.552);\draw[black!14,line width=0.3pt] (3.52,0) -- (3.52,3.552);\draw[black!14,line width=0.3pt] (3.74,0) -- (3.74,3.552);\draw[black!14,line width=0.3pt] (3.96,0) -- (3.96,3.552);\draw[black!14,line width=0.3pt] (4.18,0) -- (4.18,3.552);\draw[black!14,line width=0.3pt] (4.4,0) -- (4.4,3.552);\draw[black!14,line width=0.3pt] (4.62,0) -- (4.62,3.552);\draw[black!14,line width=0.3pt] (4.84,0) -- (4.84,3.552);\draw[black!14,line width=0.3pt] (5.06,0) -- (5.06,3.552);\draw[black!14,line width=0.3pt] (5.28,0) -- (5.28,3.552);\draw[black!14,line width=0.3pt] (5.5,0) -- (5.5,3.552);\draw[black!14,line width=0.3pt] (5.72,0) -- (5.72,3.552);\draw[black!14,line width=0.3pt] (5.94,0) -- (5.94,3.552);\draw[black!14,line width=0.3pt] (6.16,0) -- (6.16,3.552);\draw[black!14,line width=0.3pt] (6.38,0) -- (6.38,3.552);\draw[black!14,line width=0.3pt] (6.6,0) -- (6.6,3.552);\draw[black!40,line width=0.3pt] (0,0) -- (0,3.552);\node[below,inner sep=2pt] at (0,0) {0};\draw[black!40,line width=0.3pt] (1.1,0) -- (1.1,3.552);\node[below,inner sep=2pt] at (1.1,0) {50};\draw[black!40,line width=0.3pt] (2.2,0) -- (2.2,3.552);\node[below,inner sep=2pt] at (2.2,0) {100};\draw[black!40,line width=0.3pt] (3.3,0) -- (3.3,3.552);\node[below,inner sep=2pt] at (3.3,0) {150};\draw[black!40,line width=0.3pt] (4.4,0) -- (4.4,3.552);\node[below,inner sep=2pt] at (4.4,0) {200};\draw[black!40,line width=0.3pt] (5.5,0) -- (5.5,3.552);\node[below,inner sep=2pt] at (5.5,0) {250};\draw[black!40,line width=0.3pt] (6.6,0) -- (6.6,3.552);\node[below,inner sep=2pt] at (6.6,0) {300};\filldraw[fill=black!30,draw=black] (0,0.196) rectangle (3.74,0.476);\node[left,inner sep=3pt] at (0,0.336) {Mo};\filldraw[fill=black!30,draw=black] (0,0.676) rectangle (2.86,0.956);\node[left,inner sep=3pt] at (0,0.816) {Di};\filldraw[fill=black!30,draw=black] (0,1.156) rectangle (4.62,1.436);\node[left,inner sep=3pt] at (0,1.296) {Mi};\node[left,inner sep=3pt] at (0,1.776) {Do};\filldraw[fill=black!30,draw=black] (0,2.116) rectangle (4.18,2.396);\node[left,inner sep=3pt] at (0,2.256) {Fr};\filldraw[fill=black!30,draw=black] (0,2.596) rectangle (5.28,2.876);\node[left,inner sep=3pt] at (0,2.736) {Sa};\filldraw[fill=black!30,draw=black] (0,3.076) rectangle (6.16,3.356);\node[left,inner sep=3pt] at (0,3.216) {So};\draw[-{Stealth[length=2mm]}] (0,0) -- (7,0);\draw (0,0) -- (0,3.552);\node[below,inner sep=2pt,yshift=-12pt] at (3.3,0) {Anzahl der Besucher};\node[rotate=90,above,inner sep=14pt] at (0,1.776) {Wochentag};\node[above,font=\small\bfseries] at (3.3,3.602) {Besucher im Schwimmbad};\end{tikzpicture}}{In dem Diagramm soll die Anzahl der Besucher eines Schwimmbades in der ersten Ferienwoche der Sommerferien dargestellt werden. Am Donnerstag waren 180 Besucher im Schwimmbad. Ergänze im Diagramm den Balken für Donnerstag.}{0}
\fusshilfe{\gfuss{14}{Balken Do bis 180}}
\end{pfaufg}\vspace{-6pt}
% DRUCK 2021-OS-K5b deckt 2021-OS-K5b
\begin{pfaufg}{nach~P10\\’21}{15.}{}
\gtextkaro{\begin{tabular}[t]{@{}c@{}}{\footnotesize\setlength{\tabcolsep}{4pt}\renewcommand{\arraystretch}{1.2}\begin{tabular}{|l|c|}\hline \textbf{Güter} & \textbf{Angaben in \%}\\\hline Schrott & 17\\\hline Flüssigkeiten & \\\hline Container & 66\\\hline Getreide & 6\\\hline Sonstiges & 1\\\hline \textbf{Gesamt} & \textbf{100}\\\hline\end{tabular}}\\[2mm]\begin{tikzpicture}[font=\footnotesize,line width=0.6pt]\filldraw[fill=black!8,draw=black,line width=0.5pt] (0,0) -- (90:1.5) arc (90:54:1.5) -- cycle;\filldraw[fill=black!55,draw=black,line width=0.5pt] (0,0) -- (54:1.5) arc (54:32.4:1.5) -- cycle;\filldraw[fill=black!75,draw=black,line width=0.5pt] (0,0) -- (32.4:1.5) arc (32.4:-28.8:1.5) -- cycle;\filldraw[fill=black!40,draw=black,line width=0.5pt] (0,0) -- (-28.8:1.5) arc (-28.8:-32.4:1.5) -- cycle;\filldraw[fill=black!22,draw=black,line width=0.5pt] (0,0) -- (-32.4:1.5) arc (-32.4:-270:1.5) -- cycle;\draw (0,0) circle (1.5);\draw[line width=0.4pt] (0.278,0.856) -- (2,1.75) -- (4,1.75);\draw[line width=0.4pt] (0.656,0.616) -- (2,0.95) -- (4,0.95);\draw[line width=0.4pt] (0.9,0.028) -- (2,-0.35) -- (4,-0.35);\draw[line width=0.4pt] (0.775,-0.458) -- (2,-1.25) -- (4,-1.25);\draw[line width=0.4pt] (-0.789,-0.434) -- (-2,1) -- (-4,1);\end{tikzpicture}\end{tabular}}{Die Tabelle zeigt die prozentualen Anteile der gesamten Beladungen und Entladungen im Hamburger Hafen. Ergänze in der Tabelle den fehlenden Prozentsatz für „Flüssigkeiten“. Beschrifte das Kreisdiagramm mit den entsprechenden Gütern.}{3}
\fusshilfe{\gfuss{15}{10 \%; Flüssigkeiten, Getreide, Schrott, Sonstiges, Container}}
\end{pfaufg}\vspace{-6pt}
% DRUCK 2018-OS-K3c deckt 2018-OS-K3c
\begin{pfaufg}{nach~P10\\’18}{16.}{}
\gtextkaro{\begin{tabular}[t]{@{}c@{}}\begin{tikzpicture}[font=\footnotesize,line width=0.6pt]\draw[black!14,line width=0.3pt] (0,0) -- (7.875,0);\draw[black!14,line width=0.3pt] (0,0.433) -- (7.875,0.433);\draw[black!14,line width=0.3pt] (0,0.867) -- (7.875,0.867);\draw[black!14,line width=0.3pt] (0,1.3) -- (7.875,1.3);\draw[black!14,line width=0.3pt] (0,1.733) -- (7.875,1.733);\draw[black!14,line width=0.3pt] (0,2.167) -- (7.875,2.167);\draw[black!14,line width=0.3pt] (0,2.6) -- (7.875,2.6);\draw[black!40,line width=0.3pt] (0,0) -- (7.875,0);\draw[black!40,line width=0.3pt] (0,0.433) -- (7.875,0.433);\draw[black!40,line width=0.3pt] (0,0.867) -- (7.875,0.867);\draw[black!40,line width=0.3pt] (0,1.3) -- (7.875,1.3);\draw[black!40,line width=0.3pt] (0,1.733) -- (7.875,1.733);\draw[black!40,line width=0.3pt] (0,2.167) -- (7.875,2.167);\draw[black!40,line width=0.3pt] (0,2.6) -- (7.875,2.6);\draw (-0.08,0) -- (0.08,0); \node[left,inner sep=2pt] at (0,0) {0\,\%};\draw (-0.08,0.433) -- (0.08,0.433); \node[left,inner sep=2pt] at (0,0.433) {10\,\%};\draw (-0.08,0.867) -- (0.08,0.867); \node[left,inner sep=2pt] at (0,0.867) {20\,\%};\draw (-0.08,1.3) -- (0.08,1.3); \node[left,inner sep=2pt] at (0,1.3) {30\,\%};\draw (-0.08,1.733) -- (0.08,1.733); \node[left,inner sep=2pt] at (0,1.733) {40\,\%};\draw (-0.08,2.167) -- (0.08,2.167); \node[left,inner sep=2pt] at (0,2.167) {50\,\%};\draw (-0.08,2.6) -- (0.08,2.6); \node[left,inner sep=2pt] at (0,2.6) {60\,\%};\filldraw[fill=black!45,draw=black] (0.963,0) rectangle (1.663,0.433);\node[below,inner sep=3pt,align=center,text width=1.6cm] at (1.312,0) {nie};\filldraw[fill=black!45,draw=black] (2.712,0) rectangle (3.413,0.65);\node[below,inner sep=3pt,align=center,text width=1.6cm] at (3.062,0) {1-mal bis 3-mal pro Woche};\filldraw[fill=black!45,draw=black] (4.463,0) rectangle (5.162,2.383);\node[below,inner sep=3pt,align=center,text width=1.6cm] at (4.812,0) {4-mal bis 6-mal pro Woche};\filldraw[fill=black!45,draw=black] (6.213,0) rectangle (6.912,0.867);\node[below,inner sep=3pt,align=center,text width=1.6cm] at (6.562,0) {täglich};\draw[-{Stealth[length=2mm]}] (0,0) -- (0,3.05) node[above right,inner sep=1pt] {};\draw[-{Stealth[length=2mm]}] (0,0) -- (8.175,0);\node[above,font=\small\bfseries] at (3.938,3.35) {So oft essen Frauen Fleisch oder Wurst};\end{tikzpicture}\\[2mm]\begin{tikzpicture}\draw[black!16,line width=0.3pt,step=5mm] (0,0) grid (15,3);\draw[line width=1pt] (1,0.5) rectangle (11,2);\end{tikzpicture}\end{tabular}}{Frauen wurden nach ihrem Essverhalten befragt. Zeichne zu dem dargestellten Essverhalten der Frauen ein Streifendiagramm.}{0}
\fusshilfe{\gfuss{16}{nie 1 cm, 1- bis 3-mal 1,5 cm, 4- bis 6-mal 5,5 cm, täglich 2 cm}}
\end{pfaufg}\vspace{-6pt}
% DRUCK 2017-OS-K2c deckt 2017-OS-K2c
\begin{pfaufg}{nach~P10\\’17}{17.}{}
\gtextkaro{\begin{tabular}[t]{@{}c@{}}{\footnotesize\setlength{\tabcolsep}{4pt}\renewcommand{\arraystretch}{1.2}\begin{tabular}{|l|r|}\hline \textbf{Im Jahr 2013 geboren als} & \begin{tabular}{@{}c@{}}\textbf{Anteil in}\\\textbf{Prozent (\%)}\end{tabular}\\\hline erstes Kind der Mutter & 49,4\\\hline zweites Kind der Mutter & 34,4\\\hline drittes Kind der Mutter & 11,2\\\hline \begin{tabular}{@{}l@{}}viertes oder weiteres\\Kind der Mutter\end{tabular} & 5,0\\\hline\end{tabular}}\\[2mm]\begin{tikzpicture}[font=\footnotesize,line width=0.6pt]\filldraw[fill=black!30,draw=black] (0,0) -- (180:1.8) arc (180:198:1.8) -- cycle;\draw (0,0) circle (1.8);\draw[-{Stealth[length=2mm]},line width=0.5pt] (-2.5,-1.35) -- (-1.45,-0.35);\foreach \y in {-1.6,-2.15,-2.7}{\draw[line width=0.4pt,dotted] (-4.4,\y) -- (-2.4,\y);}\end{tikzpicture}\end{tabular}}{Im Jahr 2013 wurden in Deutschland 682\,069 Kinder geboren. Im Kreisdiagramm ist ein Anteil aus der Tabelle farbig dargestellt. Beschrifte diesen farbigen Anteil. Berechne die Winkelgröße, die der Anteil für „drittes Kind der Mutter“ haben muss. Zeichne den Anteil für „drittes Kind der Mutter“ in das Kreisdiagramm ein und beschrifte diesen Anteil.}{3}
\par\medskip\noindent\hfill Winkel $=$\linie{45mm}
\fusshilfe{\gfuss{17}{viertes oder weiteres Kind; 40,32°}}
\end{pfaufg}\vspace{-6pt}
% DRUCK 2015-OS-K7c deckt 2015-OS-K7c
\begin{pfaufg}{nach~P10\\’15}{18.}{}
\gtextkaro{\begin{tikzpicture}[line width=0.6pt]\draw (0,0) circle (1.8);\fill (0,0) circle (1.4pt);\draw (0,0) -- (1.8,0);\end{tikzpicture}}{In der Saison 2012\slash{}13 hat Hertha BSC in der 2.~Bundesliga gespielt. In allen Spielen der 2.~Bundesliga wurden in dieser Saison 83 Elfmeter gegeben. Davon wurden 71 verwandelt. Berechne den prozentualen Anteil der nicht verwandelten Elfmeter in dieser Saison. Zeichne diesen Anteil in das vorbereitete Kreisdiagramm ein. Gib die Größe des zugehörigen Winkels an.}{5}
\par\medskip\noindent\hfill Anteil $=$\linie{25mm}\quad Winkel $=$\linie{25mm}
\fusshilfe{\gfuss{18}{14,5 \%; 52,0°}}
\end{pfaufg}\vspace{-6pt}
\end{document}